% Part: methods
% Chapter: proofs
% Section: inference-patterns

\documentclass[../../../include/open-logic-section]{subfiles}

\begin{document}

\olfileid[id]{mth}{prf}{pat}

\olsection{Pola Inferensi}

Bukti tersusun atas inferensi-inferensi individual. Ketika membuat
inferensi, kita biasanya menandainya dengan menggunakan kata seperti
``jadi'', ``dengan demikian'', atau ``oleh karena itu''. Inferensi itu
sering bertumpu pada satu atau dua fakta yang sudah tersedia dalam
bukti---mungkin sesuatu yang telah kita asumsikan, atau sesuatu yang
telah kita simpulkan melalui inferensi sebelumnya. Agar jelas, kita
dapat memberi label pada hal-hal tersebut, lalu dalam inferensi kita
menunjukkan pernyataan lain mana yang sedang digunakan. Inferensi juga
sering memuat penjelasan tentang \emph{mengapa} kesimpulan baru kita
mengikuti hal-hal yang mendahuluinya. Ada beberapa pola inferensi umum
yang sangat sering digunakan dalam bukti; beberapa di antaranya akan
kita bahas di bawah. Pola inferensi tertentu, seperti pembuktian dengan
induksi, lebih rumit (dan akan dibahas kemudian).

Kita telah membahas satu pola inferensi: menguraikan, atau menerapkan,
suatu definisi. Ketika menguraikan definisi, kita hanya menyatakan
kembali sesuatu yang melibatkan definiendum dengan menggunakan
definiens. Misalnya, andaikan dalam suatu bukti kita telah menetapkan
bahwa $D = E$ (a). Kita dapat menerapkan definisi $=$ untuk himpunan dan
menarik inferensi: ``Dengan demikian, menurut definisi dan berdasarkan
(a), setiap !!{element} dari~$D$ merupakan !!a{element} dari~$E$, dan
sebaliknya.''

Agak membingungkan, kita sering tidak menuliskan pembenaran suatu
inferensi pada saat inferensi itu dibuat, tetapi sebelumnya. Andaikan
kita belum membuktikan bahwa $D = E$, tetapi ingin membuktikannya. Jika
$D = E$ merupakan kesimpulan yang dituju, kita juga dapat menyatakan
kembali tujuan tersebut dengan menerapkan definisi: untuk membuktikan
$D = E$, kita harus membuktikan bahwa setiap !!{element} dari~$D$
merupakan !!a{element} dari~$E$, dan sebaliknya. Jadi bukti kita akan
berbentuk: (a) buktikan bahwa setiap !!{element} dari~$D$ merupakan
!!a{element} dari~$E$; (b) setiap !!{element} dari~$E$ merupakan
!!a{element} dari~$D$; (c) oleh karena itu, berdasarkan (a) dan (b)
serta definisi $=$, $D = E$. Namun, biasanya kita tidak menulisnya
dengan cara ini. Sebagai gantinya, kita mungkin menulis seperti
berikut:
\begin{quote}
Kita ingin menunjukkan $D = E$. Menurut definisi~$=$, ini sama dengan
menunjukkan bahwa setiap !!{element} dari~$D$ merupakan !!a{element}
dari~$E$, dan sebaliknya.

(a) \dots{} (bukti bahwa setiap !!{element} dari~$D$ merupakan
!!a{element} dari~$E$) \dots

(b) \dots{} (bukti bahwa setiap !!{element} dari~$E$ merupakan
!!a{element} dari~$D$) \dots
\end{quote}

\subsection{Menggunakan Konjungsi}

Pola inferensi yang mungkin paling sederhana ialah menarik salah satu
konjungsi dari suatu konjungsi sebagai kesimpulan. Dengan kata lain:
jika kita telah mengasumsikan atau membuktikan bahwa $p$ dan~$q$, kita
berhak menyimpulkan bahwa~$p$ (dan juga bahwa~$q$). Inferensi ini begitu
mendasar sehingga sering tidak disebutkan. Misalnya, setelah
menguraikan definisi $D = E$, kita telah menetapkan bahwa setiap
!!{element} dari~$D$ merupakan !!a{element} dari~$E$, dan sebaliknya.
Dari sini kita dapat menyimpulkan bahwa setiap !!{element} dari~$E$
merupakan !!a{element} dari~$D$ (itulah bagian ``dan sebaliknya'').

\subsection{Membuktikan Konjungsi}

Kadang-kadang, apa yang diminta untuk Anda buktikan berbentuk
konjungsi; Anda akan diminta untuk ``membuktikan $p$ dan $q$.'' Dalam
hal ini, Anda hanya perlu melakukan dua hal: membuktikan $p$, lalu
membuktikan $q$. Anda dapat membagi bukti menjadi dua bagian dan, agar
jelas, memberi label pada keduanya. Ketika membuat catatan awal, Anda
mungkin menulis ``(1) Buktikan $p$'' di bagian atas halaman dan ``(2)
Buktikan $q$'' di tengah halaman. (Tentu saja, Anda mungkin tidak
secara eksplisit diminta membuktikan suatu konjungsi, tetapi mendapati
bahwa bukti Anda mengharuskan pembuktian konjungsi. Misalnya, jika
diminta membuktikan bahwa $D = E$, setelah menguraikan definisi~$=$ Anda
akan mendapati bahwa Anda harus membuktikan: setiap !!{element} dari
~$D$ merupakan !!a{element} dari~$E$ \emph{dan} setiap !!{element}
dari~$E$ merupakan !!a{element} dari~$D$.)

\subsection{Membuktikan Disjungsi}

Ketika apa yang Anda buktikan berbentuk disjungsi (yaitu, pernyataan
berbentuk ``$p$ atau $q$''), cukup ditunjukkan bahwa salah satu
disjungnya benar. Namun, pada dasarnya hampir tidak pernah terjadi
bahwa salah satu disjung langsung mengikuti asumsi teorema Anda. Yang
lebih sering terjadi, asumsi-asumsi teorema itu sendiri berbentuk
disjungsi, atau Anda sedang menunjukkan bahwa semua objek dari suatu
jenis memiliki salah satu dari dua sifat, tetapi sebagian objek
memiliki sifat pertama dan yang lain memiliki sifat kedua. Di sinilah
pembuktian dengan kasus berguna (lihat di bawah).

\subsection{Bukti Kondisional}

Banyak teorema yang akan Anda jumpai berbentuk kondisional (yaitu,
menunjukkan bahwa jika $p$ berlaku, maka $q$ juga benar). Kasus-kasus
ini mudah disiapkan---cukup asumsikan anteseden kondisional (dalam
kasus ini, $p$), kemudian buktikan kesimpulan~$q$ darinya. Jadi, jika
teorema Anda berbunyi ``Jika $p$, maka $q$'', mulailah bukti dengan
``asumsikan $p$'', dan pada akhirnya Anda seharusnya telah membuktikan
~$q$.

Kondisional dapat dinyatakan dengan berbagai cara. Jadi, sebagai ganti
``Jika $p$, maka $q$'', sebuah teorema dapat menyatakan bahwa ``$p$
hanya jika $q$'', ``$q$ jika $p$'', atau ``$q$, asalkan $p$.'' Semua
ini bermakna sama dan mengharuskan kita mengasumsikan $p$ lalu
membuktikan~$q$ dari asumsi tersebut. Ingat bahwa bikondisional (``$p$
jika dan hanya jika $q$'') sebenarnya merupakan dua kondisional yang
digabungkan: jika $p$, maka $q$, dan jika $q$, maka~$p$. Jadi, yang
perlu Anda lakukan adalah dua kali menjalankan bukti kondisional: satu
untuk kondisional pertama dan satu lagi untuk kondisional kedua.
Namun, kadang-kadang pernyataan ``jika dan hanya jika'' dapat
dibuktikan dengan merangkai sejumlah pernyataan ``jika dan hanya jika''
lain sehingga Anda berawal dari ``$p$'' dan berakhir pada ``$q$''---
tetapi dalam hal itu Anda harus memastikan bahwa setiap langkah
memang benar-benar berupa ``jika dan hanya jika''.

\subsection{Klaim Universal}

Menggunakan klaim universal itu sederhana: jika sesuatu benar untuk
apa pun, maka hal itu benar untuk setiap objek tertentu. Jadi,
misalnya, jika hipotesis bukti Anda adalah $A \subseteq B$, ini berarti
(setelah definisi~$\subseteq$ diuraikan) bahwa untuk setiap $x \in A$,
$x \in B$. Dengan demikian, jika Anda sudah mengetahui bahwa $z \in
A$, Anda dapat menyimpulkan $z \in B$.

Membuktikan klaim universal mungkin tampak sedikit rumit. Biasanya
pernyataan semacam ini berbentuk: ``Jika $x$ memiliki~$P$, maka ia
memiliki~$Q$'' atau ``Semua $P$ adalah $Q$.'' Tentu saja, bentuknya
mungkin tidak persis demikian, dan diperlukan sedikit latihan untuk
mengetahui secara tepat apa yang diminta untuk Anda buktikan. Namun,
kita sering harus membuktikan bahwa semua objek yang memiliki suatu
sifat juga memiliki sifat tertentu lainnya.

Cara membuktikan klaim universal adalah memperkenalkan nama atau
variabel bagi objek-objek yang memiliki sifat pertama, kemudian
menunjukkan bahwa objek-objek tersebut juga memiliki sifat kedua. Kita
dapat mengungkapkannya dengan mengatakan bahwa untuk membuktikan
sesuatu bagi \emph{semua}~$P$, Anda harus membuktikannya bagi
suatu~$P$ \emph{sembarang}. Nama yang diperkenalkan merupakan
nama bagi suatu~$P$ sembarang. Kita biasanya menggunakan huruf tunggal
sebagai nama bagi objek sembarang tersebut, dan huruf-huruf itu lazim
mengikuti konvensi: misalnya, $n$ digunakan untuk bilangan asli, $!A$
untuk !!{formula}, $A$ untuk himpunan, $f$ untuk fungsi, dan seterusnya.

Kuncinya ialah mempertahankan keumuman di sepanjang bukti. Anda mulai
dengan mengasumsikan bahwa suatu objek sembarang (``$x$'') memiliki
sifat~$P$, lalu menunjukkan (hanya berdasarkan definisi atau hal yang
boleh diasumsikan) bahwa $x$ memiliki sifat~$Q$. Karena Anda tidak
menetapkan secara khusus apa itu $x$, selain bahwa ia memiliki sifat
$P$, Anda dapat menyatakan bahwa segala sesuatu yang memiliki $P$
juga memiliki sifat~$Q$. Singkatnya, $x$ mewakili \emph{semua} objek
yang memiliki sifat~$P$.

\begin{prop}
  Untuk semua himpunan $A$ dan $B$, $A \subseteq A \cup B$.
\end{prop}

\begin{proof}
  Misalkan $A$ dan $B$ adalah himpunan sembarang. Kita ingin menunjukkan
  bahwa $A \subseteq A \cup B$. Menurut definisi $\subseteq$, ini sama
  dengan: untuk setiap $x$, jika $x \in A$, maka $x \in A \cup B$.
  Jadi, misalkan $x \in A$ adalah !!{element} sembarang dari~$A$. Kita
  harus menunjukkan bahwa $x \in A \cup B$. Karena $x \in A$, maka
  $x \in A$ atau $x \in B$. Dengan demikian, $x \in
  \Setabs{x}{x \in A \lor x \in B}$. Namun, menurut definisi $\cup $,
  hal itu berarti $x \in A \cup B$.
\end{proof}


\subsection{Pembuktian dengan Kasus}

Andaikan Anda memiliki suatu disjungsi sebagai asumsi atau sebagai
kesimpulan yang telah ditetapkan---Anda telah mengasumsikan atau
membuktikan bahwa $p$ atau $q$ benar. Anda ingin membuktikan $r$. Anda
melakukannya dalam dua langkah: pertama, Anda mengasumsikan bahwa $p$
benar dan membuktikan~$r$; kemudian Anda mengasumsikan bahwa $q$ benar
dan sekali lagi membuktikan~$r$. Cara ini berhasil karena kita
mengasumsikan atau mengetahui bahwa salah satu dari kedua alternatif
itu berlaku. Kedua langkah tersebut menetapkan bahwa masing-masing
alternatif sudah cukup untuk memastikan kebenaran~$r$. (Jika keduanya
benar, kita tidak hanya memiliki satu, melainkan dua alasan mengapa
$r$~benar. Tidak perlu membuktikan secara terpisah bahwa $r$~benar
dengan mengasumsikan $p$ sekaligus~$q$.) Untuk menunjukkan apa yang
sedang kita lakukan, kita mengumumkan bahwa kita ``membedakan kasus.''
Misalnya, andaikan kita mengetahui bahwa $x \in B \cup C$. $B \cup C$
didefinisikan sebagai $\Setabs{x}{x \in B \text{ atau } x \in C}$.
Dengan kata lain, menurut definisi, $x \in B$ atau $x \in C$. Untuk
membuktikan bahwa $x \in A$ dari sini, mula-mula kita mengasumsikan
bahwa $x \in B$ dan membuktikan $x \in A$ dari asumsi tersebut;
kemudian kita mengasumsikan $x \in C$ dan sekali lagi membuktikan $x
\in A$ darinya. Di bawah asumsi itu, Anda akan menulis ``Kita bedakan
kasus'', lalu di bawahnya ``Kasus (1): $x \in B$'', dan di tengah
halaman ``Kasus (2): $x \in C$.'' Setelah itu, Anda akan melengkapi
bagian atas dan bagian bawah halaman.

Pembuktian dengan kasus sangat berguna jika apa yang Anda buktikan
sendiri berbentuk disjungsi. Berikut contoh sederhana:

\begin{prop}
Andaikan $B \subseteq D$ dan $C \subseteq E$. Maka $B \cup C \subseteq
D \cup E$.
\end{prop}

\begin{proof}
  Asumsikan (a) bahwa $B \subseteq D$ dan (b) $C \subseteq E$. Menurut
  definisi, setiap $x \in B$ juga memenuhi $\in D$ (c), dan setiap
  $x \in C$ juga memenuhi $\in E$ (d). Untuk menunjukkan bahwa $B
  \cup C \subseteq D \cup E$, kita harus menunjukkan bahwa jika $x \in
  B \cup C$, maka $x \in D \cup E$ (menurut definisi $\subseteq$). $x
  \in B \cup C$ jika dan hanya jika $x \in B$ atau $x \in C$ (menurut
  definisi~$\cup$). Demikian pula, $x \in D \cup E$ jika dan hanya jika
  $x \in D$ atau $x \in E$. Jadi, kita harus menunjukkan: untuk setiap
  $x$, jika $x \in B$ atau $x \in C$, maka $x \in D$ atau $x \in E$.

  \begin{quote}
  Sejauh ini kita baru menguraikan definisi!{} Kita telah merumuskan
  kembali proposisi tanpa $\subseteq$ dan $\cup$, dan kini tinggal
  mencoba membuktikan suatu klaim kondisional universal. Berdasarkan
  pembahasan di atas, hal ini dilakukan dengan mengasumsikan bahwa $x$
  adalah suatu objek yang bagian ``jika'' diasumsikan benar untuknya,
  lalu kita akan menunjukkan bahwa bagian ``maka'' juga benar. Dengan
  kata lain, kita akan mengasumsikan bahwa $x \in B$ atau $x \in C$
  dan menunjukkan bahwa $x \in D$ atau $x \in E$.
  \footnote{Paragraf ini hanya menjelaskan apa yang sedang kita
    lakukan---paragraf ini bukan bagian dari bukti, dan Anda tidak
    perlu menguraikan semua perincian ini ketika menuliskan bukti
    sendiri.}
  \end{quote}

  Andaikan $x \in B$ atau $x \in C$. Kita harus menunjukkan bahwa $x
  \in D$ atau $x \in E$. Kita bedakan kasus.

  Kasus 1: $x \in B$. Berdasarkan (c), $x \in D$. Dengan demikian, $x
  \in D$ atau $x \in E$. (Di sini kita telah membuat inferensi yang
  dibahas dalam subbagian sebelumnya!)

  Kasus 2: $x \in C$. Berdasarkan (d), $x \in E$. Dengan demikian, $x
  \in D$ atau $x \in E$.
 \end{proof}


\subsection{Membuktikan Klaim Eksistensi}

Ketika diminta membuktikan suatu klaim eksistensi, soalnya biasanya
berbentuk ``buktikan bahwa ada~$x$ sedemikian sehingga $\dots x
\dots$'', yaitu bahwa terdapat suatu objek yang memiliki sifat yang
dideskripsikan oleh ``$\dots x \dots$''. Dalam hal ini, Anda harus
menentukan objek yang sesuai dan menunjukkan bahwa objek itu memiliki
sifat yang dipersyaratkan. Ini terdengar langsung, tetapi bukti jenis
ini dapat cukup rumit. Biasanya bukti itu melibatkan tindakan
\emph{membangun} atau \emph{mendefinisikan} suatu objek dan membuktikan
bahwa objek yang didefinisikan demikian memiliki sifat yang
dipersyaratkan. Menemukan objek yang tepat mungkin sulit, membuktikan
bahwa objek itu memiliki sifat yang dipersyaratkan mungkin sulit, dan
kadang-kadang bahkan sulit untuk menunjukkan bahwa Anda benar-benar
berhasil mendefinisikan suatu objek!{}

Pada umumnya, Anda akan menuliskannya dengan menentukan objek tersebut,
misalnya ``misalkan $x$ adalah \dots'' (dengan \dots{} menentukan objek
yang dimaksud), mungkin membuktikan bahwa $\dots$ memang
mendeskripsikan objek yang ada, lalu melanjutkan dengan menunjukkan
bahwa $x$ memiliki sifat~$Q$. Berikut contoh sederhana.

\begin{prop}
  Andaikan $x \in B$. Maka ada~$A$ sedemikian sehingga $A \subseteq B$
  dan $A \neq \emptyset$.
\end{prop}

\begin{proof}
  Asumsikan $x \in B$. Misalkan $A = \{x\}$.
  \begin{quote}
    Di sini kita telah mendefinisikan himpunan~$A$ dengan mencacah
    !!{element}-nya. Karena kita mengasumsikan bahwa $x$ adalah suatu
    objek, dan kita selalu dapat membentuk himpunan dengan mencacah
    !!{element}-nya, di sini kita tidak perlu menunjukkan bahwa kita
    telah berhasil mendefinisikan himpunan~$A$. Namun, kita masih harus
    menunjukkan bahwa $A$ memiliki sifat-sifat yang dipersyaratkan
    oleh proposisi. Tanpa itu, buktinya belum lengkap!{}
  \end{quote}
  Karena $x \in A$, $A \neq \emptyset$.
  \begin{quote}
    Hal ini bertumpu pada definisi $A$ sebagai $\{x\}$ dan fakta yang
    jelas bahwa $x \in \{x\}$ dan $x \notin \emptyset$.
  \end{quote}
  Karena $x$ merupakan satu-satunya !!{element} dari~$\{x\}$, dan $x
  \in B$, setiap !!{element} dari~$A$ juga merupakan !!a{element}
  dari~$B$. Menurut definisi~$\subseteq$, $A \subseteq B$.
\end{proof}

\subsection{Menggunakan Klaim Eksistensi}

Andaikan Anda mengetahui bahwa suatu klaim eksistensi benar (Anda
telah membuktikannya, atau klaim itu merupakan hipotesis yang dapat
Anda gunakan), misalnya ``untuk suatu~$x$, $x \in A$'' atau ``ada $x
\in A$.'' Jika ingin menggunakannya dalam bukti, Anda dapat berpura-pura
bahwa Anda mempunyai nama untuk salah satu objek yang menurut hipotesis
Anda ada. Karena $A$ memuat setidaknya satu objek, ada objek yang dapat
dirujuk oleh nama tersebut. Tentu saja, Anda mungkin tidak dapat
memilih atau mendeskripsikannya lebih lanjut (selain bahwa objek itu
memenuhi $\in A$). Namun, untuk keperluan bukti, Anda dapat berpura-pura
telah memilihnya dan memberinya nama. Penting untuk memilih nama yang
belum pernah Anda gunakan (dan tidak muncul dalam hipotesis); jika
tidak, sesuatu dapat menjadi keliru. Dalam bukti, Anda menandai hal
ini dengan beralih dari ``untuk suatu $x$, $x \in A$'' menjadi
``Misalkan $a \in A$.'' Sekarang Anda dapat bernalar tentang~$a$,
menggunakan hipotesis lain, dan seterusnya, sampai memperoleh
kesimpulan~$p$. Jika $p$ tidak lagi menyebut~$a$, maka $p$ tidak
bergantung pada asumsi bahwa $a \in A$, dan Anda telah menunjukkan
bahwa kesimpulan itu hanya mengikuti asumsi ``untuk suatu $x$, $x \in
A$.''

\begin{prop}
Jika $A \neq \emptyset$, maka $A \cup B \neq \emptyset$.
\end{prop}

\begin{proof}
  Andaikan $A \neq \emptyset$. Jadi, untuk suatu $x$, $x \in A$.
  \begin{quote}
    Di sini mula-mula kita hanya menyatakan kembali hipotesis
    proposisi. Hipotesis ini, yaitu $A \neq \emptyset$, menyembunyikan
    klaim eksistensial yang baru Anda peroleh setelah menguraikan
    beberapa definisi. Definisi $=$ memberi tahu kita bahwa $A =
    \emptyset$ jika dan hanya jika setiap $x \in A$ juga memenuhi
    $\in \emptyset$, dan setiap $x \in \emptyset$ juga memenuhi $\in
    A$. Dengan menyangkal kedua sisi, kita memperoleh: $A \neq
    \emptyset$ jika dan hanya jika terdapat suatu $x \in A$ yang
    memenuhi $\notin \emptyset$, atau terdapat suatu $x \in
    \emptyset$ yang memenuhi $\notin A$. Karena tidak ada apa pun
    yang memenuhi $\in \emptyset$, disjung kedua tidak pernah benar,
    dan ``$x \in A$ dan $x \notin \emptyset$'' menyederhana menjadi
    $x \in A$. Jadi, $A \neq \emptyset$ jika dan hanya jika untuk suatu
    $x$, $x \in A$. Ini adalah klaim eksistensi. Sekarang kita
    menggunakan klaim eksistensi tersebut dengan memperkenalkan nama
    bagi salah satu !!{element} dari~$A$:
  \end{quote}
  Misalkan $a \in A$.
  \begin{quote}
    Sekarang kita telah memperkenalkan nama bagi salah satu objek yang
    memenuhi~$\in A$. Kita akan terus bernalar tentang~$a$, tetapi
    berhati-hati agar hanya mengasumsikan bahwa $a \in A$, dan tidak
    mengasumsikan hal lain:
  \end{quote}
  Karena $a \in A$, maka $a \in A \cup B$, menurut definisi~$\cup$.
  Jadi, untuk suatu $x$, $x \in A \cup B$, yaitu $A \cup B \neq
  \emptyset$.
  \begin{quote}
    Pada langkah terakhir itu, kita beralih dari ``$a \in A \cup B$''
    ke ``untuk suatu $x$, $x \in A \cup B$.'' Pernyataan terakhir tidak
    lagi menyebut $a$, sehingga kita tahu bahwa ``untuk suatu $x$, $x
    \in A \cup B$'' mengikuti ``untuk suatu $x$, $x \in A$'' saja.
    Namun, itu berarti bahwa $A \cup B \neq \emptyset$.
  \end{quote}
\end{proof}

Mungkin merupakan kebiasaan yang baik untuk membedakan variabel
terikat seperti ``$x$'' dari nama hipotetis seperti $a$, sebagaimana
yang telah kita lakukan. Namun, dalam praktiknya kita sering tidak
membedakannya dan hanya menggunakan $x$, seperti berikut:
\begin{quote}
Andaikan $A \neq \emptyset$, yaitu ada $x \in A$. Menurut definisi
$\cup$, $x \in A \cup B$. Jadi, $A \cup B \neq \emptyset$.
\end{quote}
Namun, ketika melakukannya, Anda harus sangat berhati-hati untuk
menggunakan $x$ dan $y$ yang berbeda bagi klaim eksistensi yang
berbeda. Misalnya, berikut ini \emph{bukan} bukti yang benar untuk
``Jika $A \neq \emptyset$ dan $B \neq \emptyset$, maka $A \cap B \neq
\emptyset$'' (yang memang tidak benar).
\begin{quote}
Andaikan $A \neq \emptyset$ dan $B \neq \emptyset$. Jadi, untuk suatu
$x$, $x \in A$, dan juga untuk suatu $x$, $x \in B$. Karena $x \in A$
dan $x \in B$, kita memperoleh $x \in A \cap B$, menurut
definisi~$\cap$. Jadi, $A \cap B \neq \emptyset$.
\end{quote}
Dapatkah Anda menemukan tempat terjadinya langkah yang keliru dan
menjelaskan mengapa hasil tersebut tidak berlaku?

\end{document}
